% 真矢量重建：Q只抽取原始三环境和新增P4--P6；中列没有Q入边。
% 精确观测计数9/18/18。中列每环境3实心+3空心，其余各环境3实心。
\begin{tikzpicture}[foundation picture]
 \path[use as bounding box] (0,4) rectangle (112,82);
 \node[foundation node,fill=white,minimum width=30mm,minimum height=8mm] (population)
  at (54,77) {环境分布$\mathcal Q$};
 \foreach \x/\title/\count/\k/\n in {
  20/原始样本/9/3/3,54/{增加$n$}/18/3/6,90/{增加$K$}/18/6/3}{
  \node at (\x,65) {\title};
  \node at (\x,58) {\count 个观测};
  \node at (\x,51) {$K=\k,\ n=\n$};
 }
 % 原始/增n方案保留相同环境身份，源于同一次环境抽样。
 \foreach \offset/\n in {5/3,39/6}{
  \foreach \e/\y in {1/43.5,2/31.5,3/19.5}{
   \draw[foundation frame,rounded corners=1.3mm,fill=white] (\offset,\y-2.8) rectangle (\offset+30,\y+2.8);
   \node at (\offset+4.5,\y) {$\mathcal P_\e$};
   \foreach \i in {1,...,\n}{
    \ifnum\i<4
     \draw[FigureStroke,fill=FigureBlue,line width=.8pt] (\offset+10+3*\i,\y) circle (.9);
    \else
     \draw[FigureBlue,fill=white,line width=1pt] (\offset+10+3*\i,\y) circle (.9);
    \fi
   }
  }
 }
 \foreach \e/\y in {1/43.5,2/36.5,3/29.5,4/22.5,5/15.5,6/8.5}{
  \ifnum\e<4
   \draw[foundation frame,rounded corners=1.3mm,fill=white] (75,\y-2.5) rectangle (105,\y+2.5);
  \else
   \draw[foundation frame,rounded corners=1.3mm,fill=FigureRedFill] (75,\y-2.5) rectangle (105,\y+2.5);
  \fi
  \node at (79.5,\y) {$\mathcal P_\e$};
  \foreach \i in {1,2,3}{
   \draw[FigureStroke,fill=FigureBlue,line width=.8pt] (85+3*\i,\y) circle (.9);
  }
 }
 % 短括号明确Q关联哪一批新抽的环境，保持所有文字之外的路由。
 \draw[FigureStroke,line width=.9pt] (4.3,46.5)--(3.5,46.5)--(3.5,16.5)--(4.3,16.5);
 \draw[foundation flow] (population.west)--(0,77)--(0,31.5)--(3.5,31.5);
 \draw[FigureStroke,line width=.9pt] (106,25.5)--(108,25.5)--(108,5.5)--(106,5.5);
 \draw[foundation flow] (population.east)--(111,77)--(111,15.5)--(108,15.5);
\end{tikzpicture}%
